\section{Tiered access and evidence reporting}
\label{sec:retrieval}
\label{sec:lanes}

A query can reach the evidence through two entry points. Catalogue-first
search selects a collection, examines shelf descriptions and opens
selected documents or passages. Pinned search starts in specified
collections and returns chunks from their segment indexes. Both entry
points carry forward the
source identifiers needed to trace the selected evidence.

Figure~\ref{fig:tiered-algorithm} distinguishes how the workflow obtains
candidate evidence from how it reports the outcome. Algorithm~\ref{alg:tiered}
adds the reporting decisions and the option to check source text after scoring.

\begin{figure}[t]
\centering
\includegraphics[width=\columnwidth]{fig_tiered_workflow_vector}
\caption{Tiered access workflow. Pinned search and catalogue navigation
supply candidate references. The workflow can read candidates directly or
use Laya to admit passages before checking source text. A failure or limit
that prevents completion yields an incomplete report with the searched
scope and any available evidence.}
\label{fig:tiered-algorithm}
\end{figure}

The chat endpoint sets timeouts and turn limits. Workflow instructions
specify the retrieval sequence, batch sizes and what the agent should
report if it exhausts its call budget.

\subsection{Tier policies and the evidence gate}
\label{sec:trace}
\label{sec:gate}
\label{sec:gate-primitive}
\label{sec:twostage}

\looseness=-1 The knowledge-base chat interface defines two policies
(Table~\ref{tab:lanes}). The search policy
starts from pinned collections and calls two-stage or vector search. Each
tool applies the scope from the single collection identifier it receives,
so a call without one searches every collection
(Section~\ref{sec:exp-e2}). The
librarian policy starts from the full catalogue, uses shelf
descriptions to choose document references and directs the agent to score
them in sequential batches before reading the survivors. The two-stage
tool uses lexical search and optional graph expansion to form a candidate
pool, applies dense ranking and returns candidates for inspection.

The MCP server also provides \path{kb_hybrid_search}. It combines vector
and catalogue candidates and, when the gate is enabled, sorts retained
documents by its admission score (the verifier score of the ordering
records). The hybrid experiment used the top 30 vector
candidates per collection at a cosine threshold of 0.35, added catalogue
candidates and capped the fused pool at twelve. Section~\ref{sec:e6}
compares fusion order and verifier-score order on historical pools.

The local Laya classifier \cite{laya} is exposed as \path{kb_laya_judge}.
In reference mode it fetches candidate text, segments it and assigns each
passage an admission score in $[0,1]$ for the evidence question. Passages
that meet the threshold (0.5 by default) are eligible for the returned
pack. The gate orders them by source and position, deduplicates them and
applies a character limit. It returns document coordinates and source
spans alongside separate fetch and scoring errors, so valid passages can
remain available after a partial failure. The score governs passage
admission. The calling workflow or hybrid tool controls candidate fusion
and document order. Figure~\ref{fig:tiered-algorithm} shows where the
workflow chooses between direct reading and scoring, and why an incomplete
review is distinct from a completed review that found no support.

\begin{table*}[t]
\centering
\small
\setlength{\tabcolsep}{5pt}
\caption{Access policies over the tiered path. Modes differ in how candidate collections and documents are selected.
The hybrid row describes the experimental composition used in the evaluation.}
\label{tab:lanes}
\begin{tabular}{@{}>{\raggedright\arraybackslash}p{0.11\textwidth}>{\raggedright\arraybackslash}p{0.19\textwidth}>{\raggedright\arraybackslash}p{0.24\textwidth}>{\raggedright\arraybackslash}p{0.16\textwidth}>{\raggedright\arraybackslash}p{0.20\textwidth}@{}}
\toprule
Policy & Tiers traversed & Retrieval affordance & Budget type & Intended terminal report\\
\midrule
Search & pinned collections, segment & Two-stage or vector search & Endpoint turn limit & Answer or scoped not-found\\
Librarian & catalogue, shelf, document, segment & Catalogue navigation and passage scoring & Workflow instructions & Answer or searched scope\\
Hybrid & catalogue and segment & Combined vector and catalogue candidates & Configuration-specific candidate cap & Candidates or judged survivors\\
\bottomrule
\end{tabular}
\end{table*}

Response transport also determines which candidates the agent can see.
A recorded October 2026 incident truncated a document listing before the
agent could inspect the whole shelf. The shelf tool now caps its response and
supplies grouped and paginated views, allowing the caller to request the
remaining entries. Section~\ref{sec:eval} reports the repair check.

\begin{algorithm}[tbp]
\caption{Tiered access with optional evidence gate}
\label{alg:tiered}
\begin{algorithmic}[1]
\Require question $q$; catalogue $\mathcal{C}$; threshold $\tau$; gate flag $g$
\State $D \gets \textsc{GatherCandidates}(q,\mathcal{C})$ \Comment{chosen policy}
\If{$g$}
  \State $S \gets \textsc{LayaJudge}(D,q,\tau)$ \Comment{admit passages}
  \State Optionally check selected source spans and retain their references
\Else
  \State $S \gets \textsc{ReadSelected}(D)$
\EndIf
\If{execution incomplete}
  \State \Return incomplete status, scope and available evidence
\ElsIf{completed review finds no support}
  \State \Return scoped not-found with blind spots
\Else
  \State \Return answer with source references
\EndIf
\end{algorithmic}
\end{algorithm}

\subsection{Reporting search outcomes}
The workflow instructs the agent to cite sources when the evidence
supports an answer. After a completed review without support, it reports
the searched scope and remaining blind spots. If a failure or timeout
prevents completion, it reports the interruption and any usable evidence.
Section~\ref{sec:eval} evaluates these outcomes in observed agent trials.
